\documentclass[11pt,a4paper]{article}
% Préambule commun aux documents de PythonIPM/documentation.
% Compilation : tectonic <fichier>.tex
% tectonic compile avec XeTeX, qui lit l'UTF-8 nativement : surtout PAS inputenc ni fontenc,
% faits pour pdfLaTeX. Ce sont eux qui décodaient mal les caractères multi-octets — « cœur »
% ressortait en « cur », et les tirets cadratins en caractères parasites.
\usepackage{lmodern}
\usepackage[french]{babel}
\usepackage{microtype}
\usepackage[a4paper,margin=2.4cm]{geometry}
\usepackage{amsmath,amssymb}
\usepackage{booktabs}
\usepackage{longtable}
\usepackage{array}
\usepackage{ragged2e}
\usepackage{xcolor}
\usepackage{tcolorbox}
\usepackage{enumitem}
\usepackage{fancyhdr}
\usepackage[hidelinks]{hyperref}

\definecolor{encadre}{HTML}{F4F6F8}
\definecolor{bordure}{HTML}{4A6FA5}
\definecolor{alerte}{HTML}{B03A2E}

% Encadré « à retenir » : la lecture non technique du passage qui précède
\newtcolorbox{retenir}[1][]{
  colback=encadre, colframe=bordure, boxrule=0.6pt, arc=2pt,
  left=8pt, right=8pt, top=6pt, bottom=6pt,
  fonttitle=\bfseries, title={#1}
}

% Encadré d'avertissement : un point ouvert, un écart assumé
\newtcolorbox{attention}[1][]{
  colback=white, colframe=alerte, boxrule=0.6pt, arc=2pt,
  left=8pt, right=8pt, top=6pt, bottom=6pt,
  % le bandeau de titre a déjà la couleur d'alerte en fond : un titre `coltitle=alerte`
  % écrirait du rouge sur du rouge. On garde le blanc par défaut, comme l'encadré « retenir ».
  fonttitle=\bfseries, title={#1}
}

\newcommand{\col}[1]{\texttt{#1}}
% \ensuremath et non \textsubscript : \ipm est employé aussi bien dans le texte que dans les
% formules, et « M\textsubscript{0} » placé en mode mathématique est une erreur — silencieuse en
% PDF, mais qui fait échouer la conversion vers Word.
\newcommand{\ipm}{\ensuremath{M_{0}}}

\setlist[itemize]{itemsep=2pt, topsep=4pt}
\setlist[description]{style=nextline, leftmargin=1.2cm, itemsep=5pt}

% La source du document apparaît en tête de chaque page. Les documents RGPH la redéfinissent
% par \renewcommand{\source}{...} AVANT \begin{document}.
\newcommand{\source}{EHCVM 2021}

\pagestyle{fancy}
\fancyhf{}
\fancyhead[L]{\small\itshape IPM Côte d'Ivoire --- \source}
\fancyhead[R]{\small\thepage}
\renewcommand{\headrulewidth}{0.3pt}

\renewcommand{\arraystretch}{1.15}


\title{\bfseries Dictionnaire des sorties\\[4pt]
\large Les classeurs \col{indices\_ipm\_ci.xlsx} et \col{indices\_ipm\_international.xlsx},
colonne par colonne}
\author{Pipeline \col{PythonIPM} --- IPM Côte d'Ivoire, EHCVM 2021}
\date{Août 2026}

\begin{document}
\maketitle

\begin{abstract}
\noindent Ce document décrit chaque feuille et chaque colonne des classeurs de restitution
produits par l'étape 05 du pipeline. Pour chaque colonne : sa définition, sa formule, son unité,
et une phrase de lecture sur un cas réel. Il s'adresse autant à celui qui doit citer un chiffre
dans un rapport qu'à celui qui doit le recalculer.
\end{abstract}

\begin{attention}[Deux classeurs de structure identique]
L'étape 05 produit \textbf{deux} classeurs, qui ne diffèrent que par le périmètre des dimensions :
\col{indices\_ipm\_ci.xlsx} porte l'\textbf{IPM national} (16 indicateurs, 4 dimensions à 0,25) et
\col{indices\_ipm\_international.xlsx} l'\textbf{IPM international} (14 indicateurs, 3 dimensions à
1/3, sans l'Emploi). Colonnes, feuilles et formules sont \emph{les mêmes} : tout ce document vaut
pour les deux. Seules changent les valeurs, et le nombre de lignes de la feuille
\textsf{Contributions} (16 contre 14). Les chiffres cités en exemple ici sont ceux de l'IPM
national ; la comparaison des deux indices figure dans la \emph{Méthodologie de calcul}.
\end{attention}

\begin{retenir}[Où trouver le chiffre que je cherche]
\begin{center}
\small
\begin{tabular}{@{}ll@{}}
\textbf{La question} & \textbf{La colonne} \\[2pt]
Combien de pauvres ? & \col{H\_incidence} \\
À quel point sont-ils pauvres ? & \col{A\_intensite} \\
Quel est l'IPM ? & \col{M0\_ipm} \\
Qui risque de basculer ? & \col{vulnerables} \\
Qui est le plus en difficulté ? & \col{pauvrete\_severe} \\
Quelle région pèse le plus ? & \col{contribution\_M0}, feuille \textsf{Région} \\
Quel indicateur pèse le plus ? & \col{contribution\_M0}, feuille \textsf{Contributions} \\
Le chiffre est-il robuste ? & \col{M0\_cv} et \col{grappes} \\
Quelle est la marge d'erreur ? & \col{M0\_ic\_bas}, \col{M0\_ic\_haut} \\
\end{tabular}
\end{center}
\end{retenir}

\tableofcontents
\bigskip

\section{Les notations}

Toutes les formules de ce document utilisent trois symboles.

\begin{description}
  \item[$n_i$ --- le poids en personnes du ménage $i$.] $n_i = \col{taille\_menage}_i \times
  \col{ponderation\_menage}_i$. C'est le nombre de personnes réelles que le ménage enquêté
  représente dans la population ivoirienne. Toutes les moyennes de ce classeur sont pondérées par
  $n_i$ : elles décrivent la \textbf{population}, jamais l'échantillon.

  \item[$c_i$ --- le score de privation du ménage $i$.] La somme pondérée de ses privations,
  entre 0 (aucune privation) et 1 (privé sur les seize indicateurs).

  \item[$k$ --- le seuil de pauvreté multidimensionnelle.] Fixé à $1/3$. Un ménage est pauvre si
  $c_i \ge k$.
\end{description}

Le score censuré $c_i(k)$ vaut $c_i$ si le ménage est pauvre, et $0$ sinon.

\section{Structure du classeur}

Neuf feuilles. Les six feuilles de désagrégation et la feuille \textsf{Ensemble} ont exactement
les mêmes colonnes ; les feuilles \textsf{Contributions} et \textsf{Ponderations} ont une
structure propre.

\begin{center}
\begin{tabular}{@{}llr@{}}
\toprule
\textbf{Feuille} & \textbf{Une ligne =} & \textbf{Lignes} \\
\midrule
\textsf{Ensemble} & la Côte d'Ivoire entière & 1 \\
\textsf{Contributions} & un indicateur & 16 \\
\textsf{Ponderations} & un indicateur : sa dimension et son poids $w_j$ & 16 \\
\textsf{Milieu de résidence} & urbain ou rural & 2 \\
\textsf{Région} & une des 33 régions & 33 \\
\textsf{Département} & un des 108 départements & 108 \\
\textsf{Sous-préfecture ou commune} & une des 442 sous-préfectures & 442 \\
\textsf{Sexe du chef de ménage} & masculin ou féminin & 2 \\
\textsf{Taille du ménage} & une classe de taille & 5 \\
\bottomrule
\end{tabular}
\end{center}

\section{Les feuilles d'indices}

\emph{(\textsf{Ensemble} et les six désagrégations.)}

\subsection{Identification de la ligne}

\begin{description}
  \item[\col{variable}] Le découpage auquel appartient la ligne : \textsf{Ensemble},
  \textsf{Milieu de résidence}, \textsf{Région}, \textsf{Département},
  \textsf{Sous-préfecture ou commune}, \textsf{Sexe du chef de ménage} ou
  \textsf{Taille du ménage}. Utile quand plusieurs feuilles sont empilées.

  \item[\col{modalite}] Le groupe décrit, en clair : \textsf{Rural}, \textsf{Bounkani},
  \textsf{Féminin}, \textsf{7-9}. Les libellés proviennent des étiquettes de valeur des fichiers
  Stata d'origine, pas d'une saisie manuelle.

  \item[\col{code}] Le code numérique de la modalité dans l'EHCVM (1 = Urbain, 23 = Bounkani\ldots),
  conservé pour pouvoir rapprocher ces résultats d'autres sources. Vaut \textbf{0} pour les
  découpages construits par le pipeline, comme les classes de taille de ménage, qui n'ont pas de
  code d'enquête.
\end{description}

\subsection{Effectifs}

\begin{description}
  \item[\col{menages}] Le nombre de ménages \emph{enquêtés} dans le groupe. C'est un effectif
  d'échantillon, \textbf{non pondéré}. Il ne sert pas à décrire la population : il sert à juger de
  la robustesse du chiffre. Une région à 120 ménages donne un $H$ nettement moins précis qu'une
  région à 900.

  \item[\col{grappes}] Le nombre de \textbf{grappes} (zones de dénombrement) dont provient le
  groupe. C'est l'effectif qui détermine réellement la précision : dans un sondage à deux degrés,
  ce sont les grappes qui sont tirées au sort, pas les ménages. Une modalité à 24 ménages issus de
  2 grappes est nettement moins précise que 24 ménages issus de 24 grappes. En dessous de deux
  grappes, aucun intervalle de confiance n'est calculable.

  \item[\col{population}] Le nombre de personnes représentées : $\sum_i n_i$ sur le groupe. Cette
  fois, c'est bien une grandeur de population. La feuille \textsf{Ensemble} affiche 29\,809\,416.

  \item[\col{part\_population}] La part du groupe dans la population totale, entre 0 et 1. Somme à
  1 sur chaque feuille de désagrégation --- un contrôle automatique le vérifie.
\end{description}

\section{Les cinq indices}

Ce sont les colonnes que l'on cite. Chacune est une proportion comprise entre 0 et 1, à multiplier
par 100 pour l'exprimer en pourcentage.

\subsection*{\col{H\_incidence} --- l'incidence}

\[
H = \frac{\sum_i n_i \, \mathbf{1}(c_i \ge k)}{\sum_i n_i}
\]

\textbf{Ce que c'est.} La part de la population qui vit dans un ménage
multidimensionnellement pauvre. C'est un simple dénombrement de \emph{têtes} : le ménage compte ou
ne compte pas, l'ampleur de ses privations n'intervient pas.

\textbf{Interprétation.} $H$ répond à la question « \emph{combien} de pauvres ? ». C'est le chiffre
le plus intuitif, et donc celui qu'il faut manier avec le plus de prudence : un ménage à $c_i =
0{,}34$ et un ménage à $c_i = 0{,}95$ y pèsent identiquement.

\textbf{Lecture.} $H = 0{,}5796$ pour l'ensemble : \textbf{58,0~\% de la population ivoirienne},
soit 17,3 millions de personnes, vit dans un ménage pauvre --- IC 95~\% [54,9~\% ; 61,0~\%].
En rural, $H = 0{,}851$.

\subsection*{\col{A\_intensite} --- l'intensité}

\[
A = \frac{\sum_i n_i \, c_i(k)}{\sum_i n_i \, \mathbf{1}(c_i \ge k)}
\]

\textbf{Ce que c'est.} Le score de privation \emph{moyen des seuls pauvres}. Le numérateur ne
somme que les scores censurés, donc uniquement ceux des pauvres ; le dénominateur ne compte que
les pauvres. Les non-pauvres n'entrent ni en haut ni en bas.

\textbf{Interprétation.} $A$ répond à « \emph{à quel point} sont-ils pauvres ? ». Comme un ménage
pauvre a par définition $c_i \ge 1/3$, $A$ ne peut jamais descendre sous $0{,}333$ : c'est une
moyenne sur une population déjà sélectionnée. Il faut donc lire $A$ par rapport à ce plancher, pas
par rapport à zéro.

\textbf{Lecture.} $A = 0{,}4940$ : les personnes pauvres cumulent en moyenne \textbf{49,4~\% des
privations pondérées possibles}, contre un minimum concevable de 33,3~\%. La variation de $A$
entre groupes est faible (de 0,415 à 0,557 selon les régions) : en Côte d'Ivoire, c'est surtout le
\emph{nombre} de pauvres qui varie d'une région à l'autre, pas la gravité de leur situation.

\subsection*{\col{M0\_ipm} --- l'indice de pauvreté multidimensionnelle}

\[
\ipm = H \times A = \frac{\sum_i n_i \, c_i(k)}{\sum_i n_i}
\]

\textbf{Ce que c'est.} Le produit de l'incidence et de l'intensité --- ou, de façon équivalente,
la moyenne du score censuré sur \emph{toute} la population. Les deux écritures donnent le même
nombre ; le pipeline calcule les deux et vérifie leur égalité à chaque exécution.

\textbf{Interprétation.} \ipm{} est la « part des privations totales que subit effectivement la
population pauvre ». Sa qualité première est de bouger quand $H$ ne bouge pas : si une politique
améliore la situation de ménages qui restent pauvres, $H$ est inchangé mais $A$ baisse, donc \ipm{}
baisse. C'est pour cela que c'est \emph{lui}, et non $H$, qui constitue l'indice.

\textbf{Lecture.} $\ipm = 0{,}2863$. Un \ipm{} n'a pas d'unité parlante en soi : il se lit par
comparaison. Entre 0,058 à Abidjan et 0,481 dans le Bounkani, le rapport est de un à huit.

\subsection*{\col{vulnerables} --- la population vulnérable}

\[
\text{vulnerables} = \frac{\sum_i n_i \, \mathbf{1}\!\left(0{,}2 \le c_i < 1/3\right)}{\sum_i n_i}
\]

\textbf{Ce que c'est.} La part de la population dont le score est \emph{en dessous} du seuil de
pauvreté, mais pas de beaucoup. Ces ménages ne sont pas comptés dans $H$, ni dans $A$, ni dans
\ipm{} : la censure les a remis à zéro.

\textbf{Interprétation.} C'est une mesure de \emph{risque}, pas de pauvreté. Elle indique combien
de personnes basculeraient dans la pauvreté multidimensionnelle si une seule privation
supplémentaire survenait --- une école qui ferme, une source qui tarit. Un pays où $H$ baisse
pendant que la vulnérabilité augmente n'a pas nécessairement progressé : il a peut-être seulement
déplacé des ménages juste sous le seuil.

\textbf{Lecture.} 22,8~\% pour l'ensemble. À noter, un effet mécanique intéressant : la
vulnérabilité est plus forte là où la pauvreté est plus faible (33,5~\% en urbain contre 11,0~\%
en rural), simplement parce qu'en rural la plupart des ménages ont déjà franchi le seuil.

\subsection*{\col{pauvrete\_severe} --- la pauvreté sévère}

\[
\text{pauvrete\_severe} = \frac{\sum_i n_i \, \mathbf{1}(c_i \ge 0{,}5)}{\sum_i n_i}
\]

\textbf{Ce que c'est.} La part de la population dont le score atteint la moitié des privations
pondérées possibles.

\textbf{Interprétation.} Contrairement à la vulnérabilité, la pauvreté sévère est un
\textbf{sous-ensemble} de $H$ : tout ménage sévèrement pauvre est pauvre. Le rapport
$\text{pauvrete\_severe} / H$ indique la profondeur du problème : plus il est proche de 1, plus les
pauvres sont concentrés dans le bas de la distribution.

\textbf{Lecture.} 26,1~\% de la population, soit 45~\% des pauvres. En rural, 44,7~\% de la
population est sévèrement pauvre contre 9,4~\% en urbain --- un écart de un à cinq, bien plus
marqué que celui de $H$.

\begin{retenir}[Comment se recoupent les cinq indices]
Les trois statuts sont exclusifs deux à deux pour \col{vulnerables} et $H$ (on ne peut pas être à
la fois vulnérable et pauvre), tandis que \col{pauvrete\_severe} est inclus dans $H$.
Numériquement : 58,0~\% de pauvres, dont 26,1~\% de sévèrement pauvres ; plus 22,8~\% de
vulnérables ; le reste, 19,2~\%, n'est ni pauvre ni vulnérable.
\end{retenir}

\subsection*{\col{contribution\_M0} --- la part du groupe dans l'IPM national}

\[
\text{contribution\_M0} = \frac{\ipm^{(g)} \times \text{part\_population}^{(g)}}{\ipm^{\text{national}}}
\]

\textbf{Ce que c'est.} La part de l'IPM national imputable au groupe $g$. Elle repose sur la
\emph{décomposabilité} de \ipm{} : la somme des \ipm{} de groupe, pondérée par leur part de
population, redonne exactement l'\ipm{} national. Les contributions somment donc à 1 sur chaque
feuille --- vérifié automatiquement.

\textbf{Interprétation.} À comparer systématiquement avec \col{part\_population}. Un groupe qui
contribue plus que son poids démographique concentre la pauvreté ; l'inverse signale un groupe
relativement épargné.

\textbf{Lecture.} Le rural pèse 47,3~\% de la population et \textbf{71,9~\%} de l'IPM. Abidjan
pèse 21,5~\% de la population et seulement \textbf{4,3~\%} de l'IPM. Sur la feuille
\textsf{Ensemble}, la valeur est 1 par construction.

\section{La précision : intervalles, erreur-type, CV}

Chacun des trois indices principaux est publié avec son intervalle de confiance à 95~\%. Ce sont
des estimations d'échantillon : les citer sans marge revient à leur prêter une exactitude qu'elles
n'ont pas.

\begin{description}
  \item[\col{H\_ic\_bas}, \col{H\_ic\_haut}] Les bornes de l'intervalle de $H$. Pour l'ensemble :
  [0,5490 ; 0,6102] autour de $H = 0{,}5796$. Lecture : « l'incidence est comprise entre 54,9 et
  61,0~\% ».

  \item[\col{A\_ic\_bas}, \col{A\_ic\_haut}] Les bornes de l'intervalle de $A$. Elles sont
  toujours serrées ([0,4894 ; 0,4986] pour l'ensemble) : l'intensité est une moyenne sur les
  pauvres, une population nombreuse et homogène.

  \item[\col{M0\_ic\_bas}, \col{M0\_ic\_haut}] Les bornes de l'intervalle de \ipm{} :
  [0,2701 ; 0,3026]. Ce sont ces deux colonnes qui doivent accompagner tout \ipm{} publié.

  \item[\col{M0\_erreur\_type}] L'écart-type d'échantillonnage de \ipm{} : 0,0083 au niveau
  national. L'intervalle vaut l'estimation $\pm\, t \times$ cette valeur, avec $t$ pris dans la loi
  de Student à (nombre de grappes $-$ 1) degrés de liberté.

  \item[\col{M0\_cv}] Le \textbf{coefficient de variation} : erreur-type / estimation, soit
  2,9~\% au niveau national. C'est le critère de publication : au-delà de \textbf{20~\%}, une
  estimation est jugée trop imprécise pour être publiée ou commentée seule. Aucune région n'est
  concernée ; 15 départements sur 108 et 92 sous-préfectures sur 442 le sont.
\end{description}

\begin{attention}[Deux colonnes vides ne sont pas une erreur]
Quand un domaine ne contient qu'une seule grappe --- le cas de 211 sous-préfectures --- la
variance d'échantillonnage n'est pas estimable. Les colonnes d'intervalle restent alors vides. Le
pipeline ne fabrique jamais un intervalle qu'il ne peut pas calculer.
\end{attention}

Le calcul tient compte du plan de sondage à deux degrés de l'EHCVM : la variance est agrégée au
niveau de la \emph{grappe}, et non du ménage. La formule complète figure dans la \emph{Méthodologie
de calcul}.

\section{La feuille \textsf{Contributions}}

Une ligne par indicateur : seize dans le classeur national, quatorze dans l'international. Elle répond à une question différente des précédentes :
non plus « \emph{qui} est pauvre ? » mais « \emph{de quoi} la pauvreté est-elle faite ? ».

\begin{description}
  \item[\col{dimension}, \col{indicateur}, \col{colonne}] La dimension d'appartenance, le libellé
  publiable, et le nom technique de la colonne dans les fichiers intermédiaires.

  \item[\col{poids\_w}] Le poids $w_j$ de l'indicateur dans le score. Somme à 1 sur les seize
  lignes. Il vaut 0,0625 en Éducation, 0,0833 en Santé, 0,125 en Emploi et 0,0357 en
  conditions de vie. Dans le classeur international, où les trois dimensions restantes se
  partagent le poids de l'Emploi, il vaut respectivement 1/12, 1/9 et 1/21.

  \item[\col{taux\_privation\_censure}] Le taux de privation \textbf{censuré} $CH_j$ : la part de
  la population qui est \emph{à la fois} pauvre \emph{et} privée sur cet indicateur.
  \[
  CH_j = \frac{\sum_i n_i \, g^0_{ij} \, \mathbf{1}(c_i \ge k)}{\sum_i n_i}
  \]
  À ne pas confondre avec le taux de privation brut du tableau des définitions : l'assurance
  maladie prive 90,2~\% de la population au sens brut, mais son taux censuré est de 57,3~\% ---
  la différence, ce sont les ménages non assurés qui ne sont pas pauvres par ailleurs.

  \item[\col{apport\_a\_M0}] Le produit $w_j \times CH_j$, c'est-à-dire ce que l'indicateur ajoute
  en points d'\ipm{}. La somme des seize apports vaut exactement \ipm{} $= 0{,}2863$.

  \item[\col{contribution\_M0}] $C_j = w_j \, CH_j / \ipm$, la part de l'indicateur dans l'IPM.
  Somme à 1. L'emploi agricole de subsistance en explique 16,9~\% et l'assurance maladie 16,7~\% ;
  la promiscuité, 1,1~\%.

  \item[\col{contribution\_dimension}] La somme des contributions des indicateurs de la même
  dimension, répétée sur chaque ligne de cette dimension pour permettre un tri direct.
\end{description}

\begin{attention}[Poids nominal et poids réel]
Les quatre dimensions sont équipondérées à 25~\%, mais leurs contributions réelles vont de 18,5~\%
(Emploi) à 32,3~\% (Éducation). Le poids $w_j$ fixe ce qu'un indicateur \emph{peut} apporter au
maximum ; ce qu'il apporte réellement dépend de la fréquence de la privation. À poids identique
(0,125), l'emploi agricole de subsistance apporte 16,9~\% de l'\ipm{} et le chômage 1,6~\% :
un indicateur rare ne peut pas porter une dimension, quel que soit son poids.
\end{attention}

\section{Les autres sorties}

Le classeur est la restitution ; le pipeline produit aussi les tables de travail, dans
\col{sorties/dta/} (format Stata, types conservés) et \col{sorties/csv/} (texte).

\begin{center}
\small
\begin{tabular}{@{}l>{\RaggedRight}p{8.4cm}r@{}}
\toprule
\textbf{Table} & \textbf{Contenu} & \textbf{Dimensions} \\
\midrule
\col{preconstruction\_...} & situations brutes : effectifs concernés et défavorables, codes de
modalité, aucune valeur seuillée & 12\,965 $\times$ 38 \\
\col{vecteur\_z} & les seize seuils, avec leur source et leur énoncé & 16 $\times$ 7 \\
\col{matrice\_situationnelle\_...} & la matrice de privation $g^0$ : seize colonnes 0/1 &
12\,965 $\times$ 27 \\
\col{vecteur\_w} & les seize poids et leur dimension & 16 $\times$ 6 \\
\col{matrice\_privations\_ponderees} & $w_j\,g^0_{ij}$ et le score $c_i$, sans censure &
12\,965 $\times$ 28 \\
\col{matrice\_privations\_censuree} & $w_j\,g^0_{ij}\mathbf{1}(c_i \ge k)$, le score $c_i(k)$ et
les statuts & 12\,965 $\times$ 32 \\
\col{indices\_ipm\_ci}, \col{contributions\_ipm\_ci} & les mêmes données que le classeur
national & 593 $\times$ 21 \\
\col{indices\_ipm\_international}, \col{contributions\_ipm\_international} & idem, périmètre
à trois dimensions & 593 $\times$ 21 \\
\bottomrule
\end{tabular}
\end{center}

Dans les deux matrices de ménages, chaque ligne est identifiée par la clé \col{grappe} +
\col{menage} + \col{vague}, et porte les colonnes de pondération et de désagrégation. Les suffixes
sont systématiques : \col{\_ponderee} pour la matrice pondérée non censurée,
\col{\_censuree} pour la matrice pondérée censurée.

\begin{retenir}[Trois vérifications que l'on peut refaire soi-même]
\begin{enumerate}
  \item La somme d'une ligne de \col{matrice\_privations\_ponderees} redonne \col{score}.
  \item La somme d'une ligne de \col{matrice\_privations\_censuree} redonne \col{score\_censure}.
  \item La moyenne pondérée population de \col{score\_censure} redonne \ipm{} $= 0{,}2863$.
\end{enumerate}
Ces trois égalités sont contrôlées automatiquement, mais elles restent le moyen le plus rapide de
s'assurer qu'un fichier n'a pas été altéré.
\end{retenir}

\vfill
\noindent\rule{\linewidth}{0.4pt}\\[2pt]
\small
Documents associés : \emph{Méthodologie de calcul} (la construction complète, dont le calcul des
intervalles de confiance) et \emph{Note sur la dimension Santé} (les indicateurs qui remplacent la
mortalité juvénile et la nutrition).

\end{document}
